\chapter{Test functions and Distributions}

Our main objective in this chapter is to make useful definitions to generalize the derivative of a function. For that, we will need a mechanism that can weaken the conditions for a function to have partial derivatives, by removing it entirely and passing it to another convenient function that always has a derivative in any axis. 


\section{The elements of $\testf$ and $\dist$}

\begin{definition}{Test functions}{test}
    Let $\Omega$ be a subset of $\mathbb{R}^n$. Then $\testf$ is the set of all $C^\infty$ functions defined on $\Omega$ with compact support. 
\end{definition}

An example of such function is:

\[
    f(x) =
    \begin{cases}
    \text{exp}(\frac{-a^2}{a^2 - x^2}),\left\lvert x\right\rvert < a\\
    0, \left\lvert x\right\rvert\geq a
    \end{cases}
\]

And, for any $\varphi$ on this set and any locally integrable function $f:\mathbb{R}^n\rightarrow\mathbb{R}$, using integration by parts we get:

\begin{equation}
    \int_{K}D^\alpha f\varphi = (-1)^{\left\lvert\alpha\right\rvert}\int_{K}fD^\alpha\varphi
\end{equation}

Where $K\supseteq\text{supp}\varphi$ is a compact subset of $\Omega$ and $\alpha$ is a multi-index. This will motivate our next definition, but before that, we need to quickly describe the topology of $\testf$:

\begin{tcolorbox}[colback=gray!5,colframe=gray!80!black]
    A sequence $\{\varphi_n\}_{n\in\mathbb{N}}$ in $\testf$ converges to $\varphi$ if there exists a compact set $K\subset\Omega$ such that supp$(\varphi_n)\in K$ for all $n$ and $D^\alpha\varphi_n\rightarrow D^\alpha\varphi$ for all $\alpha$. 
\end{tcolorbox}

Now we can define the most important tool used in this endeavour:

\begin{definition}{Distributions}{dist}
    A linear functional $T$ is said to be a distribution if for all sequences $\{\varphi_n\}_{n\in\mathbb{N}}$ in $\testf$ convergent to some $\varphi$, the sequence $\{T\varphi_n\}_{n\in\mathbb{N}}$ converges to $T\varphi$.
\end{definition}

We will use $\dist$ for the set of all distributions. Note that $\dist$ uses the same notation as the dual space of $\testf$ and that is no coincidence, as the distributions are precisely the dual elements of the test functions, as they are the linear functionals defined from $\testf$ to $\mathbb{R}$. Furthermore, it is worth to point out that both $\testf$ and $\dist$ are vector spaces over $\mathbb{R}$.

Now, aiming for (2), we will define the following:

\begin{definition}{Distribution given by a function}{dist_func}
    Let $f:\mathbb{R}^n\rightarrow\mathbb{R}$ be a locally integrable function, we say that the distribution given by $f$ is:

    \[T_f(\varphi) = \int_{\Omega}f\varphi\]
    
\end{definition}

It is trivial to verify that $T_f$ is in fact a distribution, because as $\varphi_n\rightarrow\varphi$, $\int_{\Omega}f\varphi_n\rightarrow\int_{\Omega}f\varphi$.

Another useful distribution is the Dirac "$\delta$-function", which is in fact a distribution given by:
\begin{definition}{The Dirac distribution}{dirac}
    For any $\varphi\in\testf$:
    \[\delta(\varphi)=\varphi(0)\]
\end{definition}
Both distributions will be essential for the next sections.

Lastly, we'll define the support of a distribution:

\begin{definition}
    The support of a distribution $T$ is the complement of the largest set in $\Omega$ in which for all test functions $T(\varphi) = 0$
\end{definition}

By this definition we can conclude that, if $T = T_f$ for some locally integrable function, then:

\[\text{supp}(T_f)= \text{supp}(f)\]

And for the Dirac $\delta$ distribution (also for all of its derivatives that will be defined):

\[\text{supp}(\delta)= \{0\}\]

\subsection{Multiplication of a distribution by a function}

From now on, we will use the notation $\cinf$ for the $C^\infty$ functions defined in $\Omega$. First, let $\{\varphi_n\}_{n\in\mathbb{N}}$ be a sequence of functions in $\testf$ converging to $\varphi$, notice that, for $\psi\in\cinf$, ${\{\psi\varphi_n\}_{n\in\mathbb{N}}}$ converges to $\varphi$. Additionally, as $\varphi_n$ has compact support, $\psi\varphi$ will also. Therefore, if $T$ is a distribution, the sequence ${\{T(\psi\varphi_n)\}_{n\in\mathbb{N}}}$ is convergent to $T(\psi\varphi)$. So:

\[\psi T(\varphi) = T(\psi\varphi)\]
defines a distribution.
If $T$ is a distribution given by a function $f$, we get that:

\[\psi T_f(\varphi) = T_f(\psi\varphi) = \int_{\Omega}f\psi\varphi = T_{f\psi}(\varphi)\]


\section{Distributional derivatives}

Having (2) in mind, let's build a definition of a distribution derivative that can satisfy that equation. Let $f$ be a smooth function in $\Omega$, so:

\[D^{\alpha}Tf(\varphi) = D^{\alpha}\int_{\Omega}f\varphi\]

Now, by the dominated convergence theorem we can push the derivative inside the integral, that is:

\[D^{\alpha}\int_{\Omega}f\varphi = \int_{\Omega}D^{\alpha}f\varphi = (-1)^{\left\lvert\alpha\right\rvert}\int_{K}fD^\alpha\varphi\]

So, we now know what we have to define $D^{\alpha}T$ as to serve our needs. We let $f$ as a smooth function just so $D^\alpha f$ has some meaning in a sense, but in general, the only condition $f$ must abide is to be locally integrable.

\begin{definition}{Derivative of a distribution}{dist_derivative}
    Let $T\in\dist$, the $\alpha$ derivative of $T$ is:
    \[D^\alpha T(\varphi) = (-1)^{\left\lvert\alpha \right\rvert}T(D^\alpha\varphi)\]
\end{definition}

We need to verify that $D^\alpha T$ is also a distribution in $\dist$:

\begin{theorem}{}{}
    For $T\in\dist$ and any multi-index $\alpha$, $D^\alpha T$ is also in $\dist$.
\end{theorem}

\begin{proof}
    Let $\{\varphi_n\}_{n\in\mathbb{N}}$ be a sequence of functions in $\testf$ converging to $\varphi$, then by the topology of $\testf$, $\{D^\alpha\varphi_n\}_{n\in\mathbb{N}}$ converges to $D^\alpha\varphi$. Additionally, have that $\{T(\varphi_n)\}_{n\in\mathbb{N}}$ converges to $T\varphi$, therefore, by \cref{def:dist_derivative}:

    \[D^\alpha T(\varphi_n) = (-1)^{\left\lvert\alpha \right\rvert}T(D^\alpha\varphi_n)\]

    Because $\{D^\alpha\varphi\}_{n\in\mathbb{N}}$ is a converging sequence, $\{T(D^\alpha\varphi_n)\}_{n\in\mathbb{N}}$ converges to $T(D^\alpha\varphi)$ and then $D^\alpha T\in\dist$.


\end{proof}
Another important thing to note is that, for a distribution given by a function $f$, the following identity is verified:

\begin{lemma}{}{prop}
    \[D^\alpha T_f(\varphi) = T_{D^\alpha f} = (-1)T_f(D^\alpha\varphi)\]
\end{lemma}

Which was already proven before. In addition to that, we will say that a locally integrable function $f$ has a distributional derivative if there exists a distribution $G$ such that for all $\varphi\in\testf$:

\[\dv{}{t}T_f(\varphi) = G(\varphi)\]

So $G$ is the distributional derivative of $f$. If $G=T_g$ for some locally integrable function $g$, then we also say that $g$ is the distributional derivative of $f$. This distributional derivative is also what is called the weak derivative of $f$, sometimes denoted as $\partial_t f$.

It should be clear that if a function has a regular derivative, its weak derivative will be the same, as the identity above will be satisfied by:

\[\dv{}{t}T_f = T_{\dv{f}{t}}\]


\begin{example}{The derivative of the Heaviside Function}{heaviside}
    The Heaviside function is defined as:

    \[
    H(x)
    \begin{cases}
        1, x \geq 0\\
        0, x < 0
    \end{cases}
    \]

    So the distributional derivative of $H$ is the $\delta$ distribuition.
\end{example}

\begin{proof}
    Let $\Omega\in\mathbb{R}$. As $H$ is a locally integrable function, let's evaluate the distribution given by $H$:

    \[T_H(\varphi) = \int_{\Omega}H(x)\varphi(x)dx\]
So, by taking the derivative of this distribution and using \cref{lem:prop} we get:
    \[\dv{}{x}T_H(\varphi) = -T_H(\dv{\varphi}{x}) = -\int_{0}^{+\infty}H(x)\dv{\varphi}{x}(x)dx = \varphi(x)\Big|_{0}^{+\infty} = \varphi(0) = \delta(\varphi)\]
Therefore, $\delta$ is the distributional derivative of the Heaviside function.

\end{proof}

\begin{example}{The derivatives of the $\delta$ distribution}{deltader}
    For $\Omega\subseteq\mathbb{R}$ and $\varphi\in\testf$, using \cref{def:dist_derivative}, we get:
    \[\delta^{(n)}(\varphi(x)) = (-1)^n\delta(\varphi^{(n)}(x)) = (-1)^n\varphi^{(n)}(x)\]
    Where $\varphi^{(n)}$ is the $n$-nht derivative of $\varphi$.
\end{example}

We'll finish this section by giving some useful properties for the derivative of a distribution:
\begin{proposition}{The derivative of a distribuition multiplied by a $\cinf$ function}{devmult}
    By \cref{def:dist_derivative}:

    \[D^\alpha(\psi T)(\varphi) = D^\alpha T(\psi\varphi) = (-1)^{\left\lvert\alpha\right\rvert}T(D^\alpha\psi\varphi) = (-1)^{\left\lvert\alpha\right\rvert}(D^\alpha\psi)T(\varphi)\]
    
\end{proposition}

\begin{theorem}{Leibniz formula for distributions}{leibniz}
    For $\psi\in\cinf$, $T\in\dist$ and $\alpha$ a multi-index:

    \[D^\alpha(\psi T) = \sum_{\beta\leq\alpha}\frac{\alpha!}{\beta!(\alpha-\beta)!}D^\beta\psi D^{\alpha-\beta}T\]

\end{theorem}

\section{Convolutions}

\subsection{Convolution of functions}

Let's first remind ourselves of the definition of a convolution between two functions. For $f$,$g\in L_1(\mathbb{R}^n)$, the convolution of $f$ and $g$ is defined as:

\[f\ast g(x) = \int_{\mathbb{R}^n} f(x-y)g(y)dy\]

It is also important to point out that convolution is a commutative and associative operation as to keep in mind what kind of operation we want to define for distributions. Also, there are some important theorems and lemmas for convolutions of functions that are useful to have at hand.

\begin{theorem}{}{}
    Let $f\in L^1(\mathbb{R}^n)$ and, for $1<p<+\infty$, $g\in L^p(\mathbb{R}^n)$. Then, $f\ast g$ is well-defined, $f\ast g\in L^p(\mathbb{R}^n)$ and:

    \[\left\lVert f\ast g\right\rVert_{L^p(\mathbb{R}^n)}\leq\left\lVert f\right\rVert_{L^1(\mathbb{R}^n)}\left\lVert g\right\rVert_{L^p(\mathbb{R}^n)}\]
    
\end{theorem}

\begin{lemma}{Young's inequality for convolutions}{}
    For $1\leq\ p,q,r\leq +\infty$ such that:

    \[\frac{1}{p} + \frac{1}{q} = 1\]

    Then, for $f\in L^p(\mathbb{R}^n)$ and $g\in L^q(\mathbb{R}^n)$, this inequality holds: 

     \[\left\lVert f\ast g\right\rVert_{L^r(\mathbb{R}^n)}\leq\left\lVert f\right\rVert_{L^p(\mathbb{R}^n)}\left\lVert g\right\rVert_{L^q(\mathbb{R}^n)}\]

\end{lemma}

\begin{theorem}{}{}
    Let $f$ and $g$ be continuous functions. If $g$ has compact support, then:

    \[\text{supp}(f\ast g)\subset\text{supp}f + \text{supp}g\]

    Where, for $A$ and $B$ sets in $\mathbb{R}^n$:

    \[A+B=\{x+y|x\in A, y\in B\}.\]
    
\end{theorem}

And lastly:

\begin{theorem}{}{}
    Let $f$ be a continuous and $g$ continuous with compact support. If one of them is $C^\infty$ then $f\ast g$ is $C^\infty$.
\end{theorem}

All the proofs for these theorems can be found in chapter 1.5 of Topics in Functional Analysis and Applications.

\subsection{Convolution of distributions}

We want to define the convolution of a distribution and a test function to be such that the following equality holds. For $T_f$ a distribution given by a function and $\varphi\in\testf$:

\begin{equation}
    (T_f\ast\varphi)(x) = \int_{\Omega}f(x-y)\varphi(y)dy
\end{equation}

So now, let us first define some notation. Let $f$ be a function on $\mathbb{R}^n$ and $y\in\mathbb{R}^n$:

\[\tau_xf(y) = f(y-x)\]
\[\check{f}(y) = f(-y)\] 

Then it follows that:

\[\check{(\tau_x{f})}(y) = f(x-y) = \tau_{-x}\check{f}(y)\]
\[\tau_x\tau_y = \tau_{x+y}\]

For a distribution $T$, we need also to define:

\[\tau_{x}T(\varphi) = T(\tau_{-x}\varphi)\text{, for all }\varphi\in\testf\]

So finally:

\begin{definition}{Convolution of a distribution and a test function}{}
    Let $\Omega\subset\mathbb{R}^n$, $T\in\dist$ and $\varphi\in\testf$. Then we define:

    \[(T\ast\varphi)(x) = T(\tau_{-x}\check{\varphi})\]

\end{definition}

It is easy to notice that this definition satisfies $(3)$. Then we can derive some properties out of this definition:

\begin{proposition}{}{}
    \begin{enumerate}
        \item For any $x\in\mathbb{R}^n$, 
        
            \[\tau_x(T\ast\varphi)=(\tau_xT)\ast\varphi=T\ast(\tau_x\varphi)\]
    
        \item For $\alpha$ a multi-index,
        
            \[D^\alpha(T\ast\varphi)=(D^\alpha T)\ast\varphi = T\ast(D^\alpha\varphi)\]

            Then, $(T\ast\varphi)\in\cinf$
        
        \item For $\psi\in\testf$, then,
        
            \[T\ast(\varphi\ast\psi)=(T\ast\varphi)\ast\psi\]

        \item Let $T\in\dist$ such that $T\ast\varphi=0$ for all $\varphi\in\testf$. Then, $T=0$
    \end{enumerate}
\end{proposition}

\begin{proof}
    \begin{enumerate}
        \item By definition: 
        
        \[\tau_x(T\ast\varphi)(y)=(T\ast\varphi)(y-x)=T(\tau_{y-x}\check{\varphi})\]
        \[=T(\tau_{-x}\tau_{y}\check{\varphi})=(\tau_xT)(\tau_{y}\check{\varphi})=(\tau_xT)\ast\varphi(y)\]
        \[=T(\tau_y\check{(\tau_x\varphi)}) = (T\ast\tau_x\varphi)(y)\]

        \item It is sufficient to prove this only for the multi-index $\alpha=(0,\dots,1,\dots,0)$, where $\alpha_i=0$ for all $i\neq k$ and $\alpha_k = 1$. Then, for $e_k$ the $k$-th standard basic vector of $\mathbb{R}^n$:
        
        \[\pdv{}{x_k}(T\ast\varphi) = \lim_{h\rightarrow0}\frac{1}{h}(T\ast\varphi)(x)-(T\ast\varphi)(x-he_k)\]
        \[=\lim_{h\rightarrow0}\frac{1}{h}(T\ast\varphi)(x)-\tau_{he_k}(T\ast\varphi)(x)\]
        \[\text{By 1. ,}v(T\ast\varphi)(x)-T\ast(\tau_{he_k}\varphi)(x) \]
        \[=\lim_{h\rightarrow0}\frac{1}{h}T\ast(\varphi-\tau_{he_k}\varphi)(x)\]
        \[=\lim_{h\rightarrow0}T\ast\bigg(\frac{\varphi-\tau_{he_k}\varphi}{h}\bigg)(x)\]
        \[=\lim_{h\rightarrow0}T\bigg(\tau_x\bigg(\frac{\varphi-\tau_{he_k}\varphi}{h}\bigg)^\vee\bigg)\]
        \[=T\bigg(\tau_x\pdv{\varphi}{x_k}\bigg)^\vee = T\ast\bigg(\pdv{\varphi}{x_k}\bigg)\]

        Then, by iteration, it follows for any multi-index.

        (To be completed\dots)
    \end{enumerate}
\end{proof}

We must now


